
\def\PUDcodigo{EME115}

\def\PUDcodacad{04506.25}

\def\PUDdisciplina{Mecanismos}

\def\PUDsemestre{S5}

\def\PUDhoras{80}

%NOVOS ---------------------------------------------------
\def\PUDhorasTeoria{60}
\def\PUDhorasPratica{20}

\def\PUDinterECA{ \hyperref[PUD:ACE011]{Física 2 (04506.13)} }

\def\PUDatuaisECA{---}

\def\PUDrecursos{Projetor multimídia; Quadro branco e pincel; Aulas em laboratório.}

%----------------------------------------------------------

\def\PUDcreditos{4}

\def\PUDprerequisito{ \hyperref[PUD:ACE011]{ACE011} }

\def\PUDpreacad{ \hyperref[PUD:ACE011]{Física 2 (04506.13)} }

\def\PUDobjetivos{Dominar o conhecimento sobre os tipos de mecanismos; Conhecer as formas de análise cinemática dos mecanismos; Ter a capacidade de desenvolver projetos mecânicos com a aplicação da teoria de mecanismos.}

\def\PUDementa{Fundamentos da Cinemática em Mecanismos; Síntese Gráfica em Mecanismos; Análise de Posições em Mecanismos; Análise de Velocidades e Acelerações em Mecanismos; Projeto de Cames; Transmissões por Engrenagens.}

\def\PUDconteudo{ & Introdução ao estudo dos mecanismos, elementos construtivos dos mecanismos, mobilidade, análise de movimentos; 
\\
& Síntese gráfica de mecanismos, caminho função e movimento, mecanismos de retorno rápido, curvas de acoplador e tipos de mecanismos;
\\
& Sistemas de coordenadas, posição e deslocamento, análise gráfica e algébrica de mecanismos, análise de posições no mecanismo biela-manivela e mecanismos com mais de quatro barras;
\\
& Análise de velocidades, análise gráfica de velocidades, centros instantâneos de velocidade, centroides, velocidade de deslizamento, soluções analíticas para análise de velocidades e velocidade de qualquer ponto de um mecanismo;
\\
& Análise de acelerações, análise gráfica de acelerações, soluções analíticas para a análise de acelerações, aceleração em qualquer ponto de um mecanismo e pulso;
\\
& Projeto de cames, terminologias para cames, diagramas, projeto e dimensionamento de cames e considerações práticas de projeto;
\\
& Transmissões por engrenagens, lei fundamental do engrenamento, nomenclaturas e terminologias, tipos de engrenagens, transmissões e rendimento em transmissões por engrenagens;
 }

\def\PUDmetodo{Aulas expositivas em que serão abordados conteúdos teóricos através da projeção de slides, desenvolvimentos no quadro e resolução de exercícios práticos e teóricos; Desenvolvimento prático de projetos mecânicos em que serão utilizados materiais recicláveis e de sucata para construção de protótipo envolvendo a teoria de mecanismos.}

\def\PUDrecursos{Projetor multimídia; Quadro branco e pincel; Laboratórios de fabricação por remoção de material, laboratório de ajustagem e laboratório de soldagem.}

\def\PUDavalia{As 60 horas da parte teórica da disciplina serão avaliadas na forma de 03 (três) avaliações escritas e teóricas e/ou na forma de trabalhos aplicados, ambos sobre os conteúdos das unidades 1 a 7. As 20 horas da parte prática da disciplina serão avaliadas na forma de protótipo de projeto mecânico de um mecanismo aplicado em soluções de engenharia. O projeto terá equivalência à 01 (uma) avaliação.}

\def\PUDbibbasica{ &  \href{www.biblioteca.ifce.edu.br/index.asp?codigo_sophia=40297}{Norton}, R. L.  Cinemática e Dinâmica dos Mecanismos.  1ª ed.  Porto Alegre, RS: AMGH, 2010.  800p.  ISBN: 9788563308191. 
\\
&  \href{www.biblioteca.ifce.edu.br/index.asp?codigo_sophia=23882}{Hibbeler}, R.  Estática - Mecânica para Engenharia.  12ª ed.  São Paulo, SP: Pearson Prentice Hall, 2011.  512p.  ISBN: 9788576058151. 
\\
&  \href{www.biblioteca.ifce.edu.br/index.asp?codigo_sophia=62872}{Hibbeler}, R. C.  Dinâmica - Mecânica para Engenharia.  12ª ed.  São Paulo, SP: Pearson Prentice Hall, 2011.  591p.  ISBN: 9788576058144. }
 %&  SANTOS, I.  F.;  Dinâmica de Sistemas Mecânicos.  1ª edição.  São Paulo.  Editora MAKRON BOOKS.  2000.  288p.  ISBN 9788534611107. }

\def\PUDbibcomplementar{ &  \href{biblioteca.ifce.edu.br/index.asp?codigo_sophia=40282}{Norton}, Robert L. Projeto de máquinas: uma abordagem integrada. 4. ed. Porto Alegre: Bookman, 2013. 1028 p. ISBN 9788582600221.
\\
&  \href{www.biblioteca.ifce.edu.br/index.asp?codigo_sophia=63128}{Juvinall}, R. C.; Marshek, K. M.  Fundamentos do Projeto de Componentes de Máquinas.  5ª ed.  Rio de Janeiro, RJ: LTC, 2016.  562p.  ISBN: 9788521630098.
\\ 
 &  \href{biblioteca.ifce.edu.br/index.asp?codigo_sophia=23902}{Collins}, Jack A. Projeto mecânico de elementos de máquinas: uma perspectiva de prevenção da falha. Rio de Janeiro: LTC, 2006. 740 p. Inclui bibliografia e indice. ISBN 8521614756.
\\
 &  \href{biblioteca.ifce.edu.br/index.asp?codigo_sophia=24029}{Melconian}, Sarkis. Elementos de máquinas. 8. ed. São Paulo: Érica, 2007. 358 p., il. Inclui bibliografia. ISBN 9788571947030.
\\
 &  \href{biblioteca.ifce.edu.br/index.asp?codigo_sophia=41147}{Halliday}, David; {Resnick}, Robert; {Walker}, Jearl. Fundamentos de física 1: mecânica. 9. ed. Rio de Janeiro: LTC, 2013. 340p. ISBN 9788521619031. }
 %&  FLORES, P.  Projecto de Mecanismos Came-seguidor.  1 edição.  Portugal.  Editora Publindustria.  2012.  190p.  ISBN 9789728953409. }
%  FLORES, P.  Projecto de Mecanismos Came-seguidor.  1 edição.  Portugal.  Editora Publindustria.  2012.  190p.  ISBN 9789728953409.

\def\PUDautor{Venceslau Xavier de Lima Filho}

\def\PUDdata{2013-06-16}

\def\PUDrevisao{2}

\def\PUDrevisor{Francisco Nélio da Costa Freitas}

\def\PUDdatarevisao{2019-05-15}

\PUDcorpo
